\documentclass[11pt]{amsart}
\usepackage[T1]{fontenc}
\usepackage{lmodern}
\usepackage[margin=1.1in]{geometry}
\usepackage[numbers,sort&compress]{natbib}
\usepackage{amssymb,amsmath,mathtools}
\usepackage{booktabs}
\usepackage[hidelinks]{hyperref}
\usepackage{xcolor}
\setlength{\emergencystretch}{1.5em}

\newtheorem{theorem}{Theorem}[section]
\newtheorem{lemma}[theorem]{Lemma}
\newtheorem{proposition}[theorem]{Proposition}
\newtheorem{corollary}[theorem]{Corollary}
\theoremstyle{definition}
\newtheorem{definition}[theorem]{Definition}
\newtheorem{construction}[theorem]{Construction}
\newtheorem{remark}[theorem]{Remark}
\newtheorem{question}[theorem]{Question}

\newcommand{\F}{\mathbb{F}}
\newcommand{\PG}{\mathrm{PG}}
\newcommand{\GL}{\mathrm{GL}}
\newcommand{\PGL}{\mathrm{PGL}}
\newcommand{\srg}{\mathrm{srg}}
\newcommand{\Aut}{\mathrm{Aut}}
\newcommand{\Tr}{\mathrm{Tr}}
\newcommand{\Cay}{\mathrm{Cay}}
\newcommand{\todo}[1]{\textcolor{red}{[TODO: #1]}}

\title[Strongly regular graphs on $768$ vertices]{Strongly regular graphs with parameters $(768,312,136,120)$\\ and a nonabelian partial difference set}
\author{Dylan Plummer}
\address{Case Western Reserve University}
\email{dylan.plummer@case.edu}
\date{Draft, \today}

\begin{document}

\begin{abstract}
We construct strongly regular graphs with parameters $(768,312,136,120)$. This parameter set and its complement $(768,455,262,280)$ are listed as open in Brouwer's table. The vertex set is three copies of $V=\F_{16}^2$. Inside each copy, adjacency is given by eight lines of the Desarguesian spread. Between two copies it is given by a set $B\subseteq V$ of size $96$ whose Fourier transform is $\pm16$ on four punctured spread lines and vanishes elsewhere. Such sets come from a signed function defined using the norm $\F_{16}\to\F_4$, and the proof that the Fourier transform takes only these values reduces to a parity property of quadratics over $\F_4$. The construction has free parameters: $1000$ random members fall into $999$ isomorphism classes. One member is a Cayley graph on the nonabelian group $C_2^4\times(C_2^4\rtimes C_3)$. This gives a $(768,312,136,120)$ partial difference set, a parameter set that no abelian group admits. The same template predicts further parameter sets, five of them open, for which suitable signed functions would give graphs.
\end{abstract}

\maketitle

%=====================================================================
\section{Introduction}

A strongly regular graph $\srg(v,k,\lambda,\mu)$ is a $k$-regular graph on $v$ vertices, neither complete nor edgeless, in which any two adjacent vertices have exactly $\lambda$ common neighbours and any two distinct non-adjacent vertices have exactly $\mu$. Brouwer maintains a table of the feasible parameter sets with $v\le1300$ and the status of each~\cite{BrouwerTable}; see also~\cite{BvM2022}. As of November~2026 the table lists all feasible parameter sets on $768$ vertices as open or ruled out, with no construction known for any of them. Among the open ones is $(768,312,136,120)$, with eigenvalues $24^{182}$ and $(-8)^{585}$, and its complement $(768,455,262,280)$.

\begin{theorem}\label{thm:main}
There exist strongly regular graphs with parameters $(768,312,136,120)$, and hence also with parameters $(768,455,262,280)$. More precisely:
\begin{enumerate}
\item Construction~\ref{con:graph} gives a $\srg(768,312,136,120)$ for every choice of three seeds (Definition~\ref{def:seed}) with pairwise disjoint supports, and such seeds exist (Proposition~\ref{prop:normseed}).
\item One member is a Cayley graph on $G=C_2^4\times(C_2^4\rtimes C_3)$, which is $\mathrm{SmallGroup}(768,1085323)$. So $G$ contains a regular $(768,312,136,120)$ partial difference set.
\item The construction gives many non-isomorphic graphs: $1000$ random members fall into $999$ isomorphism classes.
\end{enumerate}
\end{theorem}

A $(768,312,136,120)$ partial difference set cannot live in an abelian group. By a theorem of Ma~\cite{Ma1994} (see~\cite[Prop.~1.1]{DeWinterWang}), an abelian group with a cyclic Sylow $p$-subgroup and order different from $p$ contains no nontrivial partial difference set. Every group of order $768=2^8\cdot3$ has a cyclic Sylow $3$-subgroup. So Theorem~\ref{thm:main}(2) gives a partial difference set that is ``genuinely nonabelian'' in the sense of~\cite{DPSS2023,PDSS2024}. Recent systematic studies of partial difference sets in nonabelian groups exclude order $768$ for computational reasons~\cite{NelsonSwartz}; see also~\cite{SwartzTauscheck}. The closest precedent we know is Brady's $\srg(147,66,25,33)$~\cite{Brady2025}, found as a partial difference set in nonabelian groups of order $147=3\cdot7^2$ by local search. The complement of that graph has the parameters of the case $q=7$ of the template in Section~\ref{sec:template}.

The construction is in Section~\ref{sec:construction} and the seeds in Section~\ref{sec:seeds}. The proof is a $3\times3$ block computation in the character basis (Section~\ref{sec:proof}). Section~\ref{sec:cayley} describes the Cayley member. Section~\ref{sec:template} explains why $768$ is a special case of a template that also predicts several other open parameter sets.

%=====================================================================
\section{Notation}\label{sec:notation}

Let $F=\F_{16}$, and let $K=\F_4=\{0,1,\omega,\omega^2\}\subset F$. Write $\Tr\colon F\to\F_2$ for the absolute trace, $N(a)=a^5=a\cdot a^4$ for the norm $F\to K$, and $T(a)=a+a^4$ for the trace $F\to K$. Put $V=F^2$, a vector space of dimension $8$ over $\F_2$, and pair $V$ with itself by $\langle (x,y),(a,b)\rangle=\Tr(ax+by)$. For $\beta\in V$ let $\chi_\beta(v)=(-1)^{\langle v,\beta\rangle}$. For $f\colon V\to\mathbb{Z}$ let
\[
(Hf)(v)=\sum_{\beta\in V} f(\beta)\chi_\beta(v),\qquad\text{so that } H^2=256\,I.
\]

The \emph{spread lines} are the $17$ one-dimensional $F$-subspaces $L_m=\{(a,ma):a\in F\}$ for $m\in F$ and $L_\infty=\{(0,a)\}$. We write $L^*=L\setminus\{0\}$, and we identify a set of spread lines with a set of slopes in $\PG(1,16)=F\cup\{\infty\}$. The annihilator $L^\perp=\{v:\langle v,\beta\rangle=0\ \forall\beta\in L\}$ of a spread line is again a spread line, and $(L^\perp)^\perp=L$. Since $\chi_\beta$ restricted to $L^\perp$ is trivial if $\beta\in L$ and a nontrivial character otherwise,
\begin{equation}\label{eq:line}
\sum_{v\in (L^\perp)^*}\chi_\beta(v)=\begin{cases}15,&\beta\in L,\\-1,&\beta\notin L.\end{cases}
\end{equation}

For $B\subseteq V$ let $A_B$ be the $V\times V$ matrix with $(A_B)_{x,y}=1$ if $x+y\in B$ and $0$ otherwise. Each $\chi_\beta$ is an eigenvector of $A_B$, with eigenvalue $(H1_B)(\beta)=\sum_{w\in B}\chi_\beta(w)$.

%=====================================================================
\section{Seeds}\label{sec:seeds}

\begin{definition}\label{def:seed}
A \emph{seed} is a function $z\colon V\to\{0,1,-1\}$ such that
\begin{itemize}
\item the support of $z$ is $\bigcup_{L\in S}L^*$ for a set $S$ of four spread lines, the \emph{support lines} of $z$, and
\item $(Hz)(v)\in\{10,-6\}$ for every $v\in V$.
\end{itemize}
\end{definition}

\begin{lemma}\label{lem:B}
If $z$ is a seed, then $b=(Hz+6)/16$ is the indicator function of a set $B\subseteq V$ with $|B|=96$ and
\[
H1_B=16z+96\,\delta_0 .
\]
\end{lemma}

\begin{proof}
Since $Hz$ takes the values $10$ and $-6$, $b$ is $0/1$-valued. Since $H^2=256I$ and $H\delta_0=1$, we get $Hb=(256z+6\cdot256\,\delta_0)/16=16z+96\delta_0$. Finally $|B|=(Hb)(0)=96$, because $z(0)=0$.
\end{proof}

\begin{lemma}\label{lem:GL}
Let $M\in\GL(2,F)$, acting on column vectors, and let $M^{\top}$ be its transpose. If $z$ is a seed, then so is $z\circ M^{-1}$, and $H(z\circ M^{-1})=(Hz)\circ M^{\top}$.
\end{lemma}

\begin{proof}
The map $M$ permutes the spread lines. Moreover $\langle v,M\gamma\rangle=\langle M^\top v,\gamma\rangle$, since the pairing is the trace of the $F$-bilinear dot product. Hence $\sum_\beta z(M^{-1}\beta)\chi_\beta(v)=\sum_\gamma z(\gamma)\chi_\gamma(M^\top v)$.
\end{proof}

\begin{proposition}[Norm seeds]\label{prop:normseed}
Let $h\in F$, $\mu_0\in K$ and $\nu\in K^*$. Define $z=z_{h,\mu_0,\nu}$ by
\[
z\bigl(a,(h+\mu)a\bigr)=\varepsilon_\mu\,(-1)^{[N(a)=\nu]}\qquad(\mu\in K,\ a\in F^*),\qquad \varepsilon_\mu=\begin{cases}-1,&\mu=\mu_0,\\ 1,&\mu\ne\mu_0,\end{cases}
\]
and $z=0$ elsewhere. Then $z$ is a seed with support lines $\{L_m: m\in h+K\}$.
\end{proposition}

The shear $(a,b)\mapsto(a,b+ha)$ maps $z_{0,\mu_0,\nu}$ to $z_{h,\mu_0+h,\nu}$. By Lemma~\ref{lem:GL} it therefore suffices to treat $h=0$. We need two lemmas.

\begin{lemma}[Gauss periods]\label{lem:gauss}
For $c\in K^*$,
\[
\sum_{w\in F,\ N(w)=c}(-1)^{\Tr(w)}=\begin{cases}-3,&c=1,\\ \phantom{-}1,&c\in\{\omega,\omega^2\}.\end{cases}
\]
Consequently, for $\nu\in K^*$ the function $f_\nu(a)=(-1)^{[N(a)=\nu]}$ on $F^*$ satisfies
\[
\hat f_\nu(u):=\sum_{a\in F^*}f_\nu(a)(-1)^{\Tr(au)}=\begin{cases}\phantom{-}5,& u=0\text{ or }N(u)=\nu^{-1},\\-3,&\text{otherwise.}\end{cases}
\]
\end{lemma}

\begin{proof}
Write $\Tr=\mathrm{tr}\circ T$, where $\mathrm{tr}\colon K\to\F_2$ is the trace, so $\mathrm{tr}(1)=0$ and $\mathrm{tr}(\omega)=\mathrm{tr}(\omega^2)=1$. An element $w$ with $N(w)=c$ and $T(w)=s$ is a root of $X^2+sX+c$.
\begin{itemize}
\item If $s=0$, the only such $w$ is $w=c^2\in K$.
\item If $s\neq0$, a root cannot lie in $K$, because $T$ vanishes on $K$. So such $w$ exist (exactly two of them) if and only if $X^2+sX+c$ is irreducible over $K$, that is, iff $\mathrm{tr}(c/s^2)=1$, that is, iff $s\ne c^2$.
\end{itemize}
Hence the sum equals $1+2\sum_{s\in K^*\setminus\{c^2\}}(-1)^{\mathrm{tr}(s)}=1+2\bigl(-1-(-1)^{\mathrm{tr}(c)}\bigr)$. This is $-3$ for $c=1$ and $1$ otherwise.

For the second statement, $\hat f_\nu(0)=15-2\cdot5=5$. For $u\ne0$ we have $\hat f_\nu(u)=-1-2\sum_{N(a)=\nu}(-1)^{\Tr(au)}$. Substituting $w=au$ turns the last sum into the first statement with $c=\nu N(u)$.
\end{proof}

\begin{lemma}[Parity]\label{lem:parity}
For every $(x,y)\in V$ and $n\in K^*$, the number of $\mu\in K$ with $N(x+\mu y)\in\{0,n\}$ is even.
\end{lemma}

\begin{proof}
If $y=0$ the value $N(x)$ does not depend on $\mu$, so the count is $0$ or $4$. Let $y\ne0$ and $w=x/y$. Then $N(x+\mu y)=N(y)N(w+\mu)$, and we count $\mu$ with $N(w+\mu)\in N(y)^{-1}\{0,n\}$, a set of size $2$.
\begin{itemize}
\item If $w\in K$, then $N(w+\mu)=(w+\mu)^2$ is a bijection of $K$, so the count is $2$.
\item If $w\notin K$, then $N(w+\mu)=(w+\mu)(w^4+\mu)=N(w)+\mu(\mu+T(w))$ with $T(w)\ne0$. The map $\mu\mapsto\mu^2+T(w)\mu$ is $\F_2$-linear on $K$ with kernel $\{0,T(w)\}$. So every value of $N(w+\mu)$ is taken by exactly $0$ or $2$ elements $\mu$, and the count is even.\qedhere
\end{itemize}
\end{proof}

\begin{proof}[Proof of Proposition~\ref{prop:normseed}]
Take $h=0$. For $v=(x,y)$ we have $\chi_{(a,\mu a)}(v)=(-1)^{\Tr(a(x+\mu y))}$, so $(Hz)(v)=\sum_{\mu\in K}\varepsilon_\mu\hat f_\nu(x+\mu y)$. By Lemma~\ref{lem:gauss}, $\hat f_\nu(x+\mu y)=-3+8\delta_\mu$, where $\delta_\mu=1$ if $N(x+\mu y)\in\{0,\nu^{-1}\}$ and $\delta_\mu=0$ otherwise. Since $\sum_\mu\varepsilon_\mu=2$,
\[
(Hz)(v)=-6+8\Bigl(\sum_{\mu}\delta_\mu-2\delta_{\mu_0}\Bigr).
\]
By Lemma~\ref{lem:parity}, $c=\sum_\mu\delta_\mu\in\{0,2,4\}$. If $c=0$ then $\delta_{\mu_0}=0$, and if $c=4$ then $\delta_{\mu_0}=1$. In all cases $c-2\delta_{\mu_0}\in\{0,2\}$, so $(Hz)(v)\in\{-6,10\}$.
\end{proof}

\begin{remark}[All $C_5$-invariant seeds]\label{rem:allseeds}
The scalars $C_5=\ker N\le F^*$ act on $V$. A computer search over all $\binom{17}{4}$ sets of four spread lines and all $2^{12}$ sign patterns constant on $C_5$-orbits found the following.
\begin{itemize}
\item There are exactly $4080$ $C_5$-invariant seeds.
\item They form a single $\GL(2,16)$-orbit, the orbit of the norm seeds.
\item Their support lines are exactly the $340$ four-subsets of the sublines $\PG(1,4)\subset\PG(1,16)$, which are the four-sets with cross-ratio $\omega$. Each such four-set carries $12$ seeds.
\end{itemize}
Seeds that are not $C_5$-invariant also exist: if $z$ is a seed then so is $z\chi_c$ for every $c\in V$, which corresponds to translating $B$ by $c$.
\end{remark}

%=====================================================================
\section{The construction}\label{sec:construction}

\begin{construction}\label{con:graph}
Let $z_0,z_1,z_2$ be seeds whose sets of support lines $S_0,S_1,S_2$ are pairwise disjoint, and let $B_0,B_1,B_2$ be the sets of Lemma~\ref{lem:B}. Indices are taken modulo $3$. Let $T_i=S_{i-1}\cup S_i$, a set of eight spread lines, and let $C_i=\bigcup_{L\in T_i}(L^\perp)^*$, a set of size $120$. The graph $\Gamma=\Gamma(z_0,z_1,z_2)$ has vertex set $\mathbb{Z}/3\times V$ and the following edges:
\begin{itemize}
\item $(i,x)\sim(i,y)$ if $x+y\in C_i$;
\item $(i,x)\sim(i+1,y)$ if $x+y\in B_i$.
\end{itemize}
\end{construction}

Every vertex has $120+96+96=312$ neighbours. The translations $(i,x)\mapsto(i,x+c)$ are automorphisms, so $|\Aut\Gamma|\ge256$.

Proposition~\ref{prop:normseed} gives many admissible triples. For example, take three of the four additive cosets of $K$ in $F$, say $h_0+K$, $h_1+K$ and $h_2+K$, and set $z_i=z_{h_i,\mu_i,\nu_i}$ for any $\mu_i\in h_i+K$ and $\nu_i\in K^*$. More generally, one can take any three pairwise disjoint four-subsets of sublines $\PG(1,4)$ with any of the $12$ seeds on each, multiplied by characters (Remark~\ref{rem:allseeds}).

%=====================================================================
\section{Proof of Theorem~\ref{thm:main}(1)}\label{sec:proof}

\begin{theorem}\label{thm:srg}
For every admissible triple, $\Gamma(z_0,z_1,z_2)$ is a $\srg(768,312,136,120)$ with spectrum $312^1\,24^{182}\,(-8)^{585}$.
\end{theorem}

\begin{proof}
For $\beta\in V$ let $W_\beta$ be the span of the three vectors $e_i\otimes\chi_\beta$ ($i\in\mathbb Z/3$), where $e_i\otimes\chi_\beta$ is $\chi_\beta$ on fibre $i$ and zero elsewhere. Each of the nine blocks of the adjacency matrix $A$ has the form $A_X$ for some $X\subseteq V$. So $A$ preserves $W_\beta$, and the space of functions on the vertex set is the direct sum of the $256$ three-dimensional spaces $W_\beta$. On $W_\beta$, $A$ acts by the symmetric matrix $D_\beta$ with
\begin{itemize}
\item diagonal entries $(H1_{C_i})(\beta)$, and
\item off-diagonal entries $(D_\beta)_{i,i+1}=(H1_{B_i})(\beta)=16z_i(\beta)+96\delta_{\beta,0}$.
\end{itemize}

\emph{The case $\beta=0$.} $D_0$ has diagonal $120$ and off-diagonal entries $96$, with eigenvalues $312$, $24$, $24$.

\emph{The case $\beta\ne0$.} The vector $\beta$ lies on a unique spread line. By~\eqref{eq:line} and $|T_i|=8$, the diagonal entry at $i$ is $(H1_{C_i})(\beta)=16[\beta\in\bigcup T_i]-8$, which is $8$ or $-8$. Since the $S_i$ are disjoint, at most one $z_i(\beta)$ is nonzero.
\begin{itemize}
\item If $\beta$ lies on no support line, then $D_\beta=-8I$.
\item If $\beta$ lies on a line of $S_i$, then $\beta\in\bigcup T_i$ and $\beta\in\bigcup T_{i+1}$, but $\beta\notin\bigcup T_{i+2}$. So $D_\beta$ is the direct sum of $\begin{pmatrix}8&\pm16\\\pm16&8\end{pmatrix}$, with eigenvalues $24$ and $-8$, and the $1\times1$ block $(-8)$.
\end{itemize}

There are $3\cdot4\cdot15=180$ vectors $\beta$ of the second kind and $75$ of the first. So $A$ has the eigenvalue $312$ once, $24$ with multiplicity $2+180=182$, and $-8$ with multiplicity $2\cdot180+3\cdot75=585$. A connected regular graph with exactly three distinct eigenvalues is strongly regular. Here $\mu=k+rs=312-192=120$ and $\lambda=\mu+r+s=136$.
\end{proof}

%=====================================================================
\section{A Cayley member and a partial difference set}\label{sec:cayley}

Fix $c\in F\setminus K$, and let $M=\begin{pmatrix}1&0\\c&\omega\end{pmatrix}$. Then $M^3=I$, and $M$ maps the slope $m$ to $c+\omega m$. So $M$ permutes three of the four additive cosets of $K$ cyclically:
\[
K\ \longmapsto\ c+K\ \longmapsto\ c\omega^2+K\ \longmapsto\ K .
\]
Let $z_0=z_{0,\mu_0,\nu}$ and $z_i=z_0\circ M^{-i}$. These are seeds by Lemma~\ref{lem:GL}, with pairwise disjoint support lines.

\begin{proposition}\label{prop:cayley}
Let $R=(M^\top)^{-1}$. For this triple, the map $\sigma\colon(i,x)\mapsto(i+1,Rx)$ is an automorphism of $\Gamma$. Together with the translations it generates a group $G\cong V\rtimes\langle R\rangle$ of order $768$ that acts regularly on the vertices. Hence $\Gamma\cong\Cay(G,D)$, where $D$ is the neighbourhood of $(0,0)$, and $D$ is a regular $(768,312,136,120)$ partial difference set in $G$.
\end{proposition}

\begin{proof}
By Lemma~\ref{lem:GL}, $Hz_{i+1}=(Hz_i)\circ R^{-1}$, so $B_{i+1}=RB_i$. Similarly $S_{i+1}=MS_i$, hence $T_{i+1}=MT_i$. Since $\langle Rv,M\beta\rangle=\langle v,\beta\rangle$, this gives $(ML)^\perp=R(L^\perp)$ and so $C_{i+1}=RC_i$. These relations hold for all $i$ modulo $3$ because $M^3=I$. Therefore $\sigma$ preserves both kinds of edges. The group generated by $\sigma$ and the translations is $\{(i,x)\mapsto(i+j,R^jx+c)\}$, of order $3\cdot256$, and it acts regularly.
\end{proof}

The linear map $R$ has eigenvalues $1$ and $\omega^{-1}$ on $F^2$, so $V=V_1\oplus V'$ with two $R$-invariant $F$-lines, $R$ trivial on $V_1$ and fixed-point-free on $V'$. Hence $G\cong C_2^4\times(C_2^4\rtimes C_3)$. GAP identifies $G$ as $\mathrm{SmallGroup}(768,1085323)$.

For the member with $c=t$, $\mu_0=0$ and $\nu=1$, where $F=\F_2[t]/(t^4+t+1)$, the following hold.
\begin{itemize}
\item $\Aut\Gamma\cong C_2^8\rtimes C_{15}$, of order $3840$. The vertex stabiliser is $C_5$, the scalars of norm~$1$.
\item $\Aut\Gamma$ is solvable, and its subgroups of order $768$ form a single conjugacy class. They are exactly the Hall $\{2,3\}$-subgroups.
\end{itemize}
So $\Gamma$ is a Cayley graph of $G$ and of no other group. We found a second Cayley graph on the same group, in different coordinates (Section~\ref{sec:properties}). It is not isomorphic to the first, although its automorphism group and four-clique counts are the same.

%=====================================================================
\section{Properties and non-isomorphic members}\label{sec:properties}

\begin{table}[ht]
\centering
\caption{Explicit members. Each fibre is an $\Aut$-orbit or lies in one. Four-cliques are counted through a vertex of each fibre.}\label{tab:witnesses}
\begin{tabular}{lccc}
\toprule
graph & $|\Aut\Gamma|$ & vertex orbits & four-cliques per vertex\\
\midrule
Cayley, original coordinates & $3840$ & $1$ & $445040$\\
Cayley, norm coordinates ($c=t$) & $3840$ & $1$ & $445040$\\
shear, original coordinates & $1280$ & $3\times256$ & $444230,\ 443250,\ 443760$\\
shear, norm coordinates ($h=0,t,t^3$) & $1280$ & $3\times256$ & $445040$ (each)\\
random members & $256$ (typical) & $3\times256$ & varies\\
\bottomrule
\end{tabular}
\end{table}

The witnesses found first use other coordinates. With $F=\F_2[t]/(t^4+t+1)$, so that $\omega=t^2+t$, both start from a $C_5$-invariant seed on the slopes $\{0,1,t,t^3+t\}$ (listed in the repository).
\begin{itemize}
\item The original Cayley graph transports this seed by $M=\begin{pmatrix}1&0\\ t^2+1&\omega\end{pmatrix}$. This is the matrix of Section~\ref{sec:cayley} with $c=t^2+1\notin K$.
\item The original shear graph uses the shears $(a,b)\mapsto(a,b+ha)$ with $h\in\{0,\ t^2,\ t^3+t^2+1\}$.
\end{itemize}
Every $C_5$-invariant seed is a $\GL(2,16)$-image of a norm seed (Remark~\ref{rem:allseeds}), so both are instances of Construction~\ref{con:graph}.

In the three-orbit shear graph in original coordinates, the three fibres have different four-clique counts. So no automorphism moves a fibre to another fibre, and the graph is not isomorphic to either Cayley graph.

\emph{Many graphs.} We drew $1000$ random admissible triples: three pairwise disjoint four-subsets of sublines, a random seed on each, and a random character twist of $B_0$. Every resulting graph was verified on all pairs. Their bliss canonical forms give $999$ isomorphism classes; one pair of draws produced isomorphic graphs.
\begin{itemize}
\item $996$ graphs have $|\Aut\Gamma|=256$, the translations alone, and their orbits are the three fibres.
\item $2$ have $|\Aut\Gamma|=512$, with two fibres forming one orbit.
\item $2$ have $|\Aut\Gamma|=1280$.
\end{itemize}
The four explicit graphs of Table~\ref{tab:witnesses} are pairwise non-isomorphic.

%=====================================================================
\section{A template and further parameter sets}\label{sec:template}

The proof of Theorem~\ref{thm:srg} uses only the following ingredients, and these make sense for any prime power $q$ and any number $n$ of fibres.
\begin{itemize}
\item Take $n$ copies of $\F_q^2$, joined along all pairs (fibre graph $K_n$).
\item Each pair $e$ carries a set $B_e$ with $H1_{B_e}=q\,z_e+\kappa\,\delta_0$. Here $z_e$ is supported on $r$ punctured spread lines and has $|z_e|=1$ there; for odd $q$ the characters are complex, so $z_e$ is complex-valued.
\item The supports of pairs at a common fibre are disjoint.
\item Each fibre uses the annihilators of the $(n-1)r$ support lines of its pairs.
\end{itemize}
The $2\times2$ blocks then have eigenvalues $2q-(n-1)r$ and $-(n-1)r$, and so do the characters outside the supports. Two further conditions are needed: Parseval for $B_e$, namely $\kappa^2-q^2\kappa+r(q-1)q^2=0$, and the condition that the constant characters give one of these two eigenvalues, namely $(n-1)r(q-1)-\kappa\in\{2q-(n-1)r,\,-(n-1)r\}$.

For $n=3$ the solutions form two series:
\begin{itemize}
\item $q=4r-1$ with $\kappa=2rq$;
\item $(q,r)\in\{(13,3),(16,4)\}$ with $\kappa=2q(r-1)$.
\end{itemize}
Table~\ref{tab:template} lists the cases with $v\le1300$ and $n\ge3$ that are not settled by $r=1$. The table status is from~\cite{BrouwerTable}, November 2026. For $q=3$ the template is realised, with $B$ the complement of a line, and gives the Schläfli graph. For $q=7$ the parameters are those of the complement of Brady's graph~\cite{Brady2025}. We have not checked whether that graph has this structure.

\begin{table}[ht]
\centering
\caption{Parameter sets predicted by the template ($v\le1300$, beyond $r=1$). ``Open'' means listed as open in~\cite{BrouwerTable}.}\label{tab:template}
\begin{tabular}{cccccl}
\toprule
$n$ & $q$ & $r$ & $\kappa$ & parameters & status\\
\midrule
$3$ & $7$ & $2$ & $28$ & $(147,80,46,40)$ & known~\cite{Brady2025}\\
$3$ & $11$ & $3$ & $66$ & $(363,192,106,96)$ & open\\
$3$ & $13$ & $3$ & $52$ & $(507,176,70,56)$ & open\\
$3$ & $16$ & $4$ & $96$ & $(768,312,136,120)$ & this paper\\
$3$ & $19$ & $5$ & $190$ & $(1083,560,298,280)$ & open\\
$4$ & $7$ & $2$ & $28$ & $(196,120,74,72)$ & open\\
$4$ & $13$ & $3$ & $117$ & $(676,459,314,306)$ & open\\
$4$ & $16$ & $4$ & $160$ & $(1024,660,428,420)$ & known\\
\bottomrule
\end{tabular}
\end{table}

\begin{question}
Do seeds exist for $(q,r)=(11,3)$, $(13,3)$ and $(19,5)$ with $n=3$, or for $(7,2)$ and $(13,3)$ with $n=4$? For odd $q$ the natural candidates are built from multiplicative characters on the support lines, as the norm seeds are here.
\end{question}

\begin{question}
Classify all seeds on $\F_{16}^2$, not only the $C_5$-invariant ones and their character twists. How many isomorphism classes does Construction~\ref{con:graph} produce?
\end{question}

\begin{question}
Is there a $(768,312,136,120)$ partial difference set in a group of order $768$ other than $C_2^4\times(C_2^4\rtimes C_3)$? Are there graphs with these parameters that are not of the form in Construction~\ref{con:graph}?
\end{question}

%=====================================================================
\section{Computations}\label{sec:computations}

The code and graphs are in the accompanying repository \url{https://github.com/dylan-plummer/srg768}. All checks below were run with code written independently of the original search code.
\begin{itemize}
\item \textbf{Graphs checked on all pairs} (verifying $A^2=kI+\lambda A+\mu(J-I-A)$):
  \begin{itemize}
  \item the two original witnesses and the complement of the shear witness;
  \item the norm-coordinate Cayley and shear graphs;
  \item $20$ random members, and every member generated in Section~\ref{sec:properties}.
  \end{itemize}
\item \textbf{Spectrum:} the $3\times3$ blocks $D_\beta$ were also computed directly.
\item \textbf{Lemmas:} Lemmas~\ref{lem:gauss} and~\ref{lem:parity} and all $48$ norm seeds were checked by enumeration.
\item \textbf{Partial difference set:} it was checked from its file by counting all $312^2$ quotients.
\item \textbf{Groups:} $\Aut\Gamma$ was computed with bliss through python-igraph; the group identification and the subgroup classes with GAP.
\item \textbf{Remark~\ref{rem:allseeds}:} by exhaustive enumeration.
\end{itemize}

\subsection*{Use of AI}
An AI assistant was used for the search and verification programs and a first draft of the text. \todo{author to confirm and finalize this statement.}

\bibliographystyle{plainnat}
\bibliography{refs}

\end{document}
